\documentclass[10pt,a4paper]{amsart}
\usepackage[margin=19mm]{geometry}
\usepackage{amsmath,amssymb,mathtools}
\usepackage[scr=rsfs,cal=euler]{mathalfa}
\usepackage{xfrac}
\usepackage[unicode,hidelinks,bookmarks=false]{hyperref}
\newcommand{\ie}{i.e.\ }
\newcommand{\cf}{cf.\ }
\newcommand{\resp}{respectively\ }
\newcommand{\Subsection}{\subsection}
\providecommand{\nobreakdash}{\nobreak\hbox{-}}
\setlength{\emergencystretch}{2em}
\multlinegap0pt
\usepackage{amssymb}
\usepackage{tikz}
\usepackage{float}
\newtheorem{proposition}{Proposition}
\title{On graphs with equal domination and total domination numbers}
\author{Sudip Bera}
\date{Source: arXiv:2609.03512v1 (CC BY 4.0)}
\begin{document}
\maketitle

\noindent\emph{Source and licence.} On graphs with equal domination and total domination numbers. Version arXiv:2609.03512v1 (CC BY 4.0), distributed by arXiv under the Creative Commons Attribution 4.0 International licence, http://creativecommons.org/licenses/by/4.0/, as recorded in the arXiv licence metadata for this exact version and shown on the abstract page https://arxiv.org/abs/2609.03512v1. Full licence notice: \texttt{licenses/arxiv-2609.03512-CC-BY-4.0.txt}.\par\smallskip

\noindent\textit{Editorial scope.} Counted unit: the article's proposition prop:Td=2-bipartite (source0.tex lines 127--129). The article's complete original proof of it is delivered with it (source0.tex lines 131--162, one \texttt{proof} environment of 32 lines). No mathematics is written, repaired or re-derived: the statement and the proof are the article's own lines, in the article's own notation, and no retained line is substituted or re-flowed.  No further source-local context is needed: every cross-reference of every retained line resolves inside the two delivered spans.  Nothing was omitted from any retained span, so the package carries no omission record.  No retained line contains a citation, so the wrapper carries no bibliography and no bibliography entry.  No retained line contains a figure environment, an image-inclusion command or drawing code, so no image is delivered and the package declares no asset dependency.  Editorial formatting only, in addition: an AMS \texttt{amsart} preamble that restates the article's own declarations and macros used by the retained lines, namely the environments \texttt{proposition} on separate counters, together with the standard packages those lines need.  The retained environments print in this document as Proposition 1 in order of appearance, whereas the article prints the same items with the numbering of its own full document.  The complete native source of the article as distributed in the arXiv e-print for this exact version is bundled unchanged as \texttt{sources/209304/source0.tex}.
\begin{proposition}\label{prop:Td=2-bipartite}
Let $G$ be a connected graph with $\gamma(G)=2$ and $g(G)\neq 3$. If $\gamma_t(G)=2,$ then $G$ is bipartite.
\end{proposition}

\begin{proof}
Let $D=\{u,v\}$ be a total dominating set of $G$. 
Denote
\[
N(u)=\{u_1,\dots,u_k\}
\quad \text{and} \quad
N(v)=\{v_1,\dots,v_r\}.
\]
Since $D$ is minimum dominating set $N(u)$ and $N(v)$ are non-empty.
We prove the following two claims.

\medskip
\noindent
\textit{Claim I.} $N(u)$ and $N(v)$ are independent sets.

Suppose that $u_i\sim u_j$ for some $i\neq j$. Then
$u_i \sim u \sim u_j \sim u_i$ forms a cycle of length $3$, contradicting the assumption that $g(G)\neq 3$. Hence, $N(u)$ is an independent set. By a similar argument, $N(v)$ is also an independent set.

\medskip
\noindent
\textit{Claim II.} $N(u)\cap N(v)=\emptyset$.

Suppose there exists a vertex $x\in N(u)\cap N(v)$. Then $x\sim u$ and $x\sim v$. Since $u\sim v$, it follows that
$u \sim x \sim v \sim u$
is a cycle of length $3$, again contradicting $g(G)\ge 4$. Therefore, $N(u)\cap N(v)=\emptyset$. Now set
\[
A=\{u\}\cup N(v)
\quad \text{and} \quad
B=\{v\}\cup N(u).
\]
By Claims I and II, both $A$ and $B$ are independent sets. Moreover, since $D=\{u,v\}$ is a dominating set, every vertex of $G$ is adjacent to either $u$ or $v$, and hence $V(G)=A\cup B.$ Thus, $G$ admits a bipartition $(A,B)$, and therefore $G$ is bipartite.
\end{proof}

\end{document}
